\documentclass[conference]{IEEEtran}
\usepackage[utf8]{inputenc}
\usepackage[T1]{fontenc}
\usepackage{cite}
\usepackage{amsmath,amssymb}
\usepackage{graphicx}
\usepackage{booktabs}
\usepackage{url}

\title{Tool‑Augmented Post‑Processing of Identical LLM Candidates for Intent‑Driven Network Configuration (shared\_initial\_candidate evaluation)}

\author{
\IEEEauthorblockN{Autonomous AI Researcher}
\IEEEauthorblockA{CNSM 2026 Agentic AI Researcher Track}
}

\begin{document}

\maketitle

\begin{abstract}
We preregistered and executed a paired confirmatory study (cand\_netops\_toolagent\_003\_v1) to isolate whether deterministic post‑processing (topology‑aware deterministic validator + lightweight simulator) changes a binary validator pass/fail verdict when both pipelines receive the identical single LLM candidate (shared\_initial\_candidate semantics). Design: 300 paired tasks (one shared generation reused verbatim by both pipelines); primary estimand = paired\_success\_rate\_difference\_guarded\_minus\_baseline. Completed: 598 episodes \ensuremath{\rightarrow} 299 analyzable pairs. Result (artifact‑grounded): contingency n11=299, n10=0, n01=0, n00=0 \ensuremath{\rightarrow} paired diff = 0.00; exact McNemar two‑sided p = 1.0; paired bootstrap 95\% percentile CI = [0.00,0.00] (10k resamples, seed=1729). Model\_calls\_used = 300; aggregated repair\_applied=true count = 0. All raw responses, provider traces, scoring JSONs (validator\_trace included), analysis inputs/outputs, and the preregistration are archived and SHA‑256 hashed for audit (see Disclosure). We explain the observed ceiling effect, give an explicit pipeline step‑by‑step with artifact pointers proving shared reuse and repair accounting, and describe reproducible next steps to measure end‑to‑end tool‑augmentation benefits.
\end{abstract}

\section{Introduction}
Generative LLMs are being proposed to automate NetOps tasks such as intent\ensuremath{\rightarrow}configuration translation and configuration repair. Common agentic designs couple an LLM generator with deterministic validators/simulators to raise operational trustworthiness. Published claims about "tool‑augmented agents" mix generator and post‑processing effects. We therefore preregistered an isolation experiment asking: When the identical single LLM candidate is supplied verbatim to both a baseline and a guarded pipeline (shared\_initial\_candidate semantics), does deterministic post‑processing change the binary deterministic‑validator pass/fail outcome? Our contribution: a preregistered, artifact‑archived, paired confirmatory run and a concise, artifact‑grounded accounting of pipeline semantics, call counts, and validator traces supporting the observed outcome.
Research question (closed form): Under shared\_initial\_candidate semantics, does the guarded (tool‑augmented) post‑processing pipeline produce a higher deterministic‑validator pass rate than the baseline pipeline? Contributions: (1) a preregistered paired experiment (cand\_netops\_toolagent\_003\_v1); (2) an explicit pipeline workflow and per‑artifact proof of candidate reuse and repair accounting; (3) the archived confirmatory outcome (n11=299, no discordant pairs) and a discussion of ceiling effects and next steps.

\section{Related Work}
Prior work motivates tool‑augmented architectures and intent\ensuremath{\rightarrow}config tasks (examples used for positioning and verified in the evidence bundle: Nguyen et al. (2024) Intent‑Based configurations [doi:10.1002/nem.2313], PreConfig (arXiv:2403.09369), and practical tool‑augmented AIOps case studies (Rong Lu, 2026 preprint)). Surveys of LLM agents and RAG/validator designs (several reviews in the archived bibliography) argue validators increase operational safety. Our study is methodological: it isolates the post‑processing effect under a strict shared‑candidate semantics, complementing prior work that mixes generation and validation.

\section{Methodology}
Preregistration and frozen plan. The preregistration cand\_netops\_toolagent\_003\_v1 (SHA‑256 991674d5...) fixed the primary estimand = paired\_success\_rate\_difference\_guarded\_minus\_baseline, shared\_initial\_candidate semantics, sample N=300 pairs (150 simple / 150 complex), analysis executor = paired\_binary\_analysis\_v1, exclusion/missingness rules, maximum\_model\_calls = 600, and stopping rules. (Preregistration file SHA‑256 cited in Disclosure.)
Pipelines executed (artifact‑proven step‑by‑step). Under the shared\_initial\_candidate design we executed the following per task:
1) LLM generation (single call) \ensuremath{\rightarrow} produce shared\_initial candidate; response saved at execution/responses/task-<id>-shared-initial.txt (example: task-000001 SHA‑256 8f38c8be...). Provider trace for that call stored at execution/provider\_calls/task-<id>-shared-initial.json (example: task-000001 SHA‑256 fd046204...).
2a) Baseline scoring: feed the exact shared\_initial candidate (byte‑identical) into the deterministic validator (deterministic\_netops\_validator\_v5). Validator produces scoring JSON execution/scoring/task-<id>-baseline.json (contains validator\_trace, final\_state, repair\_applied flag). Representative file: execution/scoring/task-000001-baseline.json (SHA‑256 6569b491...).
2b) Guarded scoring: feed the identical shared\_initial candidate (same file) into the guarded path that also invokes the same deterministic validator and optional repair loop per the frozen retry policy. Validator produces execution/scoring/task-<id>-guarded.json. Representative file: execution/scoring/task-000001-guarded.json (SHA‑256 cefa6344...).
Note: the experiment reuses the identical candidate file for both conditions (proof in the paired response and scoring artifacts). Repair protocol and budgeting. The preregistered contract allowed up to one repair call per task; repair calls would be recorded as additional hosted\_model\_call events and produce repair\_provider\_trace files. Aggregated model\_calls\_used is recorded in the execution manifest (execution/execution\_manifest.json; model\_calls\_used = 300; SHA‑256 in manifest). Deterministic validator behavior and outputs. The validator normalizes commands, enforces encoded invariants (make‑before‑break, shutdown\ensuremath{\rightarrow}change\ensuremath{\rightarrow}restore, minimum active), produces final\_state snapshots and Violation lists, and writes validator\_trace inside each scoring JSON. Every scoring JSON therefore provides an auditable record of validator reasoning and whether any repair was applied (validator\_trace.repair\_applied boolean).

\section{Results}
Execution and basic counts (artifact pointers). Planned episodes: 600 (300 pairs). Completed episodes: 598; failed episodes: 2 (execution/execution\_manifest.json). Completed analyzable pairs: 299 (complete pairs). Hosted model call audit (execution manifest): hosted\_model\_call\_event\_count = 300; completed\_hosted\_model\_call\_event\_count = 299; failed\_hosted\_model\_call\_event\_count = 1; model\_calls\_used = 300 (one generation per task). Representative provider trace for shared call: execution/provider\_calls/task-000001-shared-initial.json (SHA‑256 fd046204...).
Primary confirmatory numbers (archived analysis artifacts). Paired contingency (analysis/paired\_contingency\_table.csv SHA‑256 d77df43a...): n11 = 299, n10 = 0, n01 = 0, n00 = 0. Baseline success rate = 299/299 = 100.0\%. Guarded success rate = 299/299 = 100.0\%. Observed paired difference = 0.00.
Statistical testing (archived analysis log). Exact McNemar two‑sided p = 1.0 (no discordant pairs). Paired bootstrap percentile 95\% CI = [0.00,0.00] (10,000 resamples; seed = 1729). Evidence artifacts: analysis/analysis\_log.jsonl (SHA‑256 7364d6a1...) and analysis/paired\_contingency\_table.csv (SHA‑256 d77df43a...).
Repair accounting (artifact grounding). Aggregated repair\_applied=true count across completed episodes = 0 (no repair provider calls observed). Evidence: model\_calls\_used = 300 (execution manifest); no repair\_provider\_trace files referenced in scoring JSONs (see representative execution/scoring/task-000001-baseline.json where validator\_trace.repair\_applied = false).

\section{Discussion}
Interpretation (artifact‑grounded): Under the strict shared\_initial\_candidate semantics and the synthetic task pool used, deterministic post‑processing (validator + simulator + optional repair loop) did not change the binary acceptance decision: every analyzable candidate accepted by the baseline path was accepted by the guarded path (n11=299). The archived scoring JSONs display validator\_trace with final\_state and repair\_applied=false for representative episodes (see execution/scoring/task-000001-baseline.json SHA‑256 6569b491... and execution/scoring/task-000001-guarded.json SHA‑256 cefa6344...).
Ceiling effect and statistical note (artifact evidence): The observed degenerate contingency (all pairs n11) produces exact McNemar p = 1.0 and a degenerate bootstrap CI [0.00,0.00]. This is the mathematical consequence of zero discordant pairs; the analysis log (analysis/analysis\_log.jsonl SHA‑256 7364d6a1...) records the bootstrap parameters. The preregistered power calculation assumed baseline \ensuremath{\approx} 0.60 and a minimally important paired improvement of 0.20; the observed baseline = 1.00 creates a ceiling effect that prevented testing that preregistered effect. This limitation and the artifact evidence are described and archived.
Operational implications and recommended follow‑ups (artifact‑referenced): (1) To quantify whether tool‑augmentation improves candidate quality researchers must permit generation to differ across conditions (no shared reuse), or inject adversarial/harder tasks that raise baseline failure probability; execution/task\_manifest.jsonl and analysis artifacts provide the format for such replications. (2) Track repair\_applied flags and repair\_provider\_trace files to measure repair loop contribution; our archived scoring JSONs show the schema and the repair\_applied field to enable reproducible accounting. (3) Replicate across model families and production CLI repertoires for external validity.

\section{Limitations}
\begin{itemize}
\item Shared\_initial semantics isolate only post‑processing and therefore cannot assess whether tool‑augmented generation (allowing generator repairs or verifier‑guided generation) would improve candidate quality.
\item Task set is synthetic and executed in a lightweight simulator; external validity to vendor‑specific CLIs, firmware idiosyncrasies, and production telemetry is limited.
\item Only one declared model family (gpt‑5‑mini) was used; results do not generalize across models without replication.
\item Deterministic validator encodes curated workflow invariants (make‑before‑break, shutdown\_configure\_restore, minimum‑active rules) but cannot exhaustively represent all production device behaviors or operator policies.
\item Observed baseline success = 100\% produced a ceiling effect, rendering the preregistered minimally important effect untestable in this run.
\end{itemize}

\section{Conclusion}
We preregistered and archived a narrow confirmatory experiment to test whether deterministic post‑processing changes validator pass/fail when identical LLM candidates are reused (shared\_initial\_candidate). For 299 analyzable complete pairs we observed a degenerate outcome: n11=299 and zero discordant pairs \ensuremath{\rightarrow} paired diff = 0.00; exact McNemar p = 1.0; bootstrap 95\% CI = [0.00,0.00]. Model\_calls\_used = 300 and no repair\_applied flags were observed in completed episodes. All artifacts (responses, provider traces, scoring JSONs with validator\_trace, task manifests, execution manifest, and analysis outputs) are archived and SHA‑256 hashed for audit. The result is narrowly scoped: it shows that when generator outputs already satisfy the validator invariants, deterministic post‑processing provides auditable traces but does not change the binary acceptance decision. Broader claims about tool‑augmented agent benefits require experiments that (i) allow per‑condition generation, (ii) inject adversarial tasks to avoid ceiling effects, and (iii) replicate across models and device repertoires. The full reproducibility bundle is available and referenced in the Disclosure Statement below for independent verification.

\section*{Disclosure Statement}
Disclosure Statement (required; archived references and SHA‑256s):
- Initial master prompt (immutable archive reference): provenance/master\_prompt.txt --- SHA‑256: 1872df1e1805d2d96940456ca016bd665d1d5196add77f5acdf1582bb39b15ba. This file is the exact frozen initial master prompt required by the CNSM Agentic AI Researcher track and documents the autonomy boundary.
- Preregistration (frozen experiment plan): cand\_netops\_toolagent\_003\_v1 --- SHA‑256: 991674d51794b41cf66471509149237b9a99f709b69aa3408d3d9d2ea2d50fc0. The preregistration fixes primary estimand, shared\_initial\_candidate semantics, N=300 paired tasks, analysis executor, and exclusion rules; it is archived unmodified.
- Declared LLM and configuration: gpt-5-mini (declared model) --- execution/model\_configuration.json (SHA‑256: 1acce92558edeb13cd97b75817329d332afe4a36d01d566f1a26c61b132923fa). Adapter family: hosted\_netops\_gvr\_v1 (execution manifest archive). All hosted provider calls are recorded in execution/provider\_calls/*.json and hashed (examples cited below).
- Deterministic validator / scorer: deterministic\_netops\_validator\_v5 (scorer\_version 5.0). Every scoring JSON written by the run contains a validator\_trace object with final\_state, normalized\_configuration, violation list, difficulty metadata, and repair\_applied flag (examples: execution/scoring/task-000001-baseline.json SHA‑256 6569b491dbfb2..., execution/scoring/task-000001-guarded.json SHA‑256 cefa6344...).
- Shared\_initial\_candidate semantics and reuse proof: each task produced a single LLM generation saved at execution/responses/task-<id>-shared-initial.txt and the identical file was supplied byte‑for‑byte to both baseline and guarded scoring paths. Representative evidence: execution/responses/task-000001-shared-initial.txt (SHA‑256 8f38c8becd7e1e36...), execution/scoring/task-000001-baseline.json (SHA‑256 6569b491...), execution/scoring/task-000001-guarded.json (SHA‑256 cefa6344...). The full set of responses and scoring JSONs are archived and hashed (see execution/execution\_manifest.json for the full per‑file manifest and SHA‑256 map).
- Repair loop accounting and model calls: The preregistered maximum\_model\_calls = 600 (one generation + at most one repair). Observed: model\_calls\_used = 300 (execution/execution\_manifest.json; SHA‑256s included in the manifest). Aggregated repair\_applied=true count across completed episodes = 0 (no repair provider traces used). Representative artifacts demonstrating no repairs applied: execution/scoring/task-000001-baseline.json where validator\_trace.repair\_applied = false (SHA‑256 6569b491...), and the provider audit entries in execution/execution\_manifest.json (hosted\_model\_call\_event\_count = 300; completed\_hosted\_model\_call\_event\_count = 299; failed\_hosted\_model\_call\_event\_count = 1). Thus the run consumed only the single shared generation per task in completed episodes.
- Missingness / failure: One provider failure caused a single pair exclusion (pair‑000269). Provider failure and logs are archived: execution/provider\_calls/task-000269-shared-initial.json (see execution manifest). The preregistered missingness rule (exclude incomplete pairs) was applied and recorded in analysis/missingness\_summary.csv (SHA‑256 0e909d00...).
- Analysis artifacts and exact numeric outputs: analysis/paired\_contingency\_table.csv (SHA‑256 d77df43abf83e...), analysis/analysis\_log.jsonl (SHA‑256 7364d6a1...), and analysis/missingness\_summary.csv (SHA‑256 0e909d00...) contain the exact numbers used in the paper (n11=299, n10=0, n01=0, n00=0; McNemar p = 1.0; bootstrap 95\% CI = [0.00,0.00], 10k resamples, seed=1729). These files are archived and hashed for audit and reproducibility.
- Human involvement and autonomy boundary: Pre‑lock human activity included broad topic selection and development/testing of the autonomous pipeline. After the autonomy lock the autonomous system performed literature discovery, preregistration, experiment design, code generation, execution, analysis, peer review, manuscript drafting, and revisions. No human manual edits of technical manuscript content were made after the autonomy lock; all deviations, restarts, and the single provider failure are logged in the execution manifest and execution\_log.jsonl. The master prompt archive documents this autonomy boundary (provenance/master\_prompt.txt SHA‑256 above).
- Representative artifact pointers (examples for quick audit):
  \textbullet{} execution/execution\_manifest.json (full per‑file manifest \& hashes; model\_calls\_used = 300)
  \textbullet{} execution/responses/task-000001-shared-initial.txt --- SHA‑256 8f38c8be...
  \textbullet{} execution/provider\_calls/task-000001-shared-initial.json --- SHA‑256 fd046204...
  \textbullet{} execution/scoring/task-000001-baseline.json --- SHA‑256 6569b491...
  \textbullet{} execution/scoring/task-000001-guarded.json --- SHA‑256 cefa6344...
  \textbullet{} analysis/paired\_contingency\_table.csv --- SHA‑256 d77df43a...
  \textbullet{} analysis/analysis\_log.jsonl --- SHA‑256 7364d6a1...
  \textbullet{} execution/model\_configuration.json --- SHA‑256 1acce925...
  \textbullet{} preregistration JSON --- SHA‑256 991674d5...
- Limitations (artifact‑grounded): (1) shared\_initial semantics isolate only post‑processing and cannot assess whether tool‑augmented generation (when generation is allowed to differ between conditions) improves candidate quality; (2) the task set is synthetic and executed in a lightweight simulator --- external validity to vendor CLIs and production devices is limited; (3) a single declared model family (gpt‑5‑mini) was used; cross‑model generality requires replication; (4) the deterministic validator encodes curated invariants but is not an exhaustive model of production device behavior; (5) observed baseline success = 100\% produced a ceiling effect preventing the preregistered minimally important effect (0.20) from being tested. All limitations and their supporting artifacts are archived and hashed.
- Availability: the full reproducibility bundle (responses, provider traces, scoring JSONs with validator\_trace, task manifest, execution manifest, analysis artifacts, preregistration, and per‑file SHA‑256 map) is archived in execution/ and analysis/ in the run bundle referenced in the submission. Inspect the listed files and SHA‑256s above for independent verification.

\begin{thebibliography}{99}
\bibitem{ref1} Nguyen Tu, Sukhyun Nam, James Won‐Ki Hong, "Intent‐Based Network Configuration Using Large Language Models", 2024, 10.1002/nem.2313
\bibitem{ref2} http://arxiv.org/abs/2403.09369
\bibitem{ref3} Rong Lu, "TADI: Tool-Augmented Drilling Intelligence via Agentic LLM Orchestration over Heterogeneous Wellsite Data", 2026, 10.20944/preprints202604.1820.v1
\bibitem{ref4} YONGGANG WENG, "WeClaw: An Open-Source Adaptive Runtime Framework for Tool-Augmented Desktop LLM Agents", 2026, 10.2139/ssrn.7317098
\bibitem{ref5} Roberto Lorusso, Antonio Maci, Antonio Coscia, "GOLLUM: Guiding cOnfiguration of firewaLL Through aUgmented Large Language Models", 2025, 10.5220/0013221900003890
\bibitem{ref6} Lingzhe Zhang, Tong Jia, Mengxi Jia, Yifan Wu, Aiwei Liu, Yong Yang, Zhonghai Wu, Xuming Hu, Philip S. Yu, Ying Li, "A Survey of AIOps in the Era of Large Language Models", 2025, 10.1145/3746635
\bibitem{ref7} Andres M. Bran, Sam Cox, Oliver Schilter, Carlo Baldassari, Andrew Dickson White, Philippe Schwaller, "Augmenting large language models with chemistry tools", 2024, 10.1038/s42256-024-00832-8
\bibitem{ref8} Wenqi Fan, Yujuan Ding, Liangbo Ning, Shijie Wang, Hengyun Li, Dawei Yin, Tat‐Seng Chua, Qing Li, "A Survey on RAG Meeting LLMs: Towards Retrieval-Augmented Large Language Models", 2024, 10.1145/3637528.3671470
\bibitem{ref9} Fuliang Li, Haozhi Lang, Jiajie Zhang, Jiaxing Shen, Xingwei Wang, "PreConfig: A Pretrained Model for Automating Network Configuration", 2024, 10.48550/arxiv.2403.09369
\bibitem{ref10} Yagmur Yigit, Mohamed Amine Ferrag, Mohamed Chahine Ghanem, Iqbal H. Sarker, Λέ\ensuremath{\alpha}ν\ensuremath{\delta}\ensuremath{\rho}ος Μ\ensuremath{\alpha}\ensuremath{\gamma}\ensuremath{\lambda}\ensuremath{\alpha}\ensuremath{\rho}άς, Christos Chrysoulas, Naghmeh Moradpoor, Norbert Tihanyi, Helge Janicke, "Generative AI and LLMs for Critical Infrastructure Protection: Evaluation Benchmarks, Agentic AI, Challenges, and Opportunities", 2025, 10.3390/s25061666
\bibitem{ref11} Wafaa Kasri, Yassine Himeur, Hamzah Ali Alkhazaleh, Saed Tarapiah, Shadi Atalla, Wathiq Mansoor, Hussain Al-Ahmad, "From Vulnerability to Defense: The Role of Large Language Models in Enhancing Cybersecurity", 2025, 10.3390/computation13020030
\bibitem{ref12} Saurabh Pahune, Zahid Akhtar, "Transitioning from MLOps to LLMOps: Navigating the Unique Challenges of Large Language Models", 2025, 10.3390/info16020087
\bibitem{ref13} Lei Wang, Chen Ma, Xueyang Feng, Zeyu Zhang, Hao Yang, Jingsen Zhang, Zhiyuan Chen, Jiakai Tang, Xu Chen, Yankai Lin, Wayne Xin Zhao, Zhewei Wei, Ji-Rong Wen, "A survey on large language model based autonomous agents", 2024, 10.1007/s11704-024-40231-1
\bibitem{ref14} Daniel Fernandes, João P. Matos-Carvalho, Carlos M. Fernandes, Nuno Fachada, "DeepSeek-V3, GPT-4, Phi-4, and LLaMA-3.3 Generate Correct Code for LoRaWAN-Related Engineering Tasks", 2025, 10.3390/electronics14071428
\bibitem{ref15} Lei Wang, Chen Ma, Xueyang Feng, Zeyu Zhang, Hao Yang, Jingsen Zhang, Zhiyuan Chen, Jiakai Tang, Xu Chen, Yankai Lin, Wayne Xin Zhao, Zhewei Wei, Ji-Rong Wen, "A survey on large language model based autonomous agents", 2024, 10.1007/s11704-024-40231-1
\end{thebibliography}

\end{document}
